%قبل از شروع مسئله ترجیحاً چیزی ننویسید. اما اگر خواستید بنویسید، در قالب لیستی (فقرات و شمارش) بنویسید.
%برای اضافه کردن نمره، بعدی از شروع‌مسئله، [x] را اضافه کنید و به جای x  نمره مورد نظر را قرار دهید. برای زیرمسئله‌ها هم به همین صورت است.
%\شروع‌مسئله[10]
\شروع‌مسئله [33]
%\متن‌سیاه{سوال عملی اول:}\\
%\متن‌سیاه{دراپ‌تست}\\
یک روز طبق معمول، عیسی و موسی در حال گذران وقت بودند، که ناگهان موسی برای سرگرم کردن عیسی فکری به ذهنش رسید. او یک درخت دودویی تقریباً کامل کشید و از عیسی خواست که کمترین تعداد تعویض خانه‌ها برای اینکه درخت، یک درخت دودویی جستجو باشد را به او بگوید (عیسی می‌تواند مقدار هر دو رأس دلخواه از درخت را با یکدیگر عوض کند). عیسی هم که تا جواب سؤال موسی را نمی‌داد، ول‌کن نبود، تصمیم گرفت برنامه‌ای بنویسد تا دیگر موسی نتواند وقت او را با این کارهای بی‌ارزش بگیرد. از آنجایی که عیسی کارهای باارزش‌تری برای انجام دارد، از شما خواسته تا این کار را برایش بکنید.
\par
در ورودی برنامه، درخت دودویی به صورت آرایه به شما داده‌می‌شود که مقدار ریشهٔ آن $a_0$ است؛ مقدار بچهٔ چپ ریشه، $a_1$، و مقدار بچهٔ راست آن $a_2$ است. به طور کلی برای رأس با مقدار $a_k$، مقدار بچهٔ چپ آن$a_{2k+1}$، و مقدار بچهٔ راست آن $a_{2k+2}$ است.

\متن‌سیاه{ورودی:}
\par
در خط اول، \متن‌لاتین{n}، تعداد رئوس درخت داده‌می‌شود.\\
سپس در خط دوم \متن‌لاتین{n} عدد که با فاصله جدا شده‌اند داده‌می‌شود ($a_0$ تا $a_{n-1}$) که مقادیر رئوس درخت هستند.
\par
\متن‌سیاه{خروجی:}
\par
در تنها سطر خروجی، حداقل تعداد تعویض‌های لازم برای تبدیل درخت به یک درخت جستجوی دودویی را چاپ نمایید.
\par
\متن‌سیاه{محدودیت‌ها:}
\par
$1 \leq n \leq 10^{5}$\\
$0 \leq a_i \leq 10^5$\par

\متن‌سیاه{ورودی و خروجی نمونه:}
\par
تست اول
\begin{latin}
	\begin{tabularx}{\textwidth}{|X | X|}
		\hline
		\textbf{stdin} &
		\textbf{stdout}\\
		\hline
		3 & 1\\
		1 2 3&\\
		\hline
	\end{tabularx}
\end{latin}
%}
% صورت اصلی مسئله را در این‌جا بنویسید.
%اگر مسئله چند قسمت داشت به این صورت عمل کنید:
%\شروع{زیرمسائل}
%\زیرمسئله[0] % مسئلۀ ۱آ
%\زیرمسئله % مسئلۀ ۱ب
%\زیرمسئله % مسئلۀ ۱پ
%\زیرمسئله % مسئلۀ ۱ت
%\پایان{زیرمسائل}

\solution{
	پاسخ نمونه.
}




%%%%%%%%%%%%%%%%%%%%%%%%%%%%%%%%%%%%%%%%%%%%%%%%%%%%%
% برخی دستورات مفید برای تمرین‌ها. %
%%%%%%%%%%%%%%%%%%%%%%%%%%%%%%%%%%%%%%%%%%%%%%%%%%%%%

% فرمول‌های نمونه:
%$f(x) = \Theta(x)$
%$T(x, y) \leq \log(x) + 2^y + \binom{2n}{n}$
% {\tt code\_function}


% لیست موارد بدون شماره‌گذاری:
%\شروع{فقرات}
%    \فقره مورد اول 
%        \شروع{فقرات}
%            \فقره مورد اول داخلی 
%                \شروع{فقرات}
%                    \فقره اولین مورد خیلی داخلی 
%                    \فقره دومین مورد خیلی داخلی 
%                \پایان{فقرات}
%            \فقره مورد دوم داخلی 
%        \پایان{فقرات}
%    \فقره مورد دوم 
%\پایان{فقرات}

% لیست موارد شماره‌گذاری شده:
%\شروع{شمارش}
%    \فقره یک مورد 
%    \فقره یک مورد دیگر 
%    \فقره باز هم یک مورد دیگر
%\پایان{شمارش}

% برای نوشتن یک فرمول چند خطی که یک شماره بخوره. برای حالت‌های دیگر به سایت زیر مراجعه کنید:
% https://www.overleaf.com/learn/latex/Aligning_equations_with_amsmath
%\begin{align} 
%    \begin{split}
%        (x+y)^3 &= (x+y)^2(x+y) \\
%                &= (x^2+2xy+y^2)(x+y) \\
%                &= (x^3+2x^2y+xy^2) + (x^2y+2xy^2+y^3) \\
%                &= x^3+3x^2y+3xy^2+y^3
%    \end{split}                                 
%\end{align}

% نوشتن ماتریس و آرایه دوبعدی
%\begin{align}
%    A = 
%    \begin{bmatrix}
%        A_{11} & A_{21} \\
%        A_{21} & A_{22}
%    \end{bmatrix}
%\end{align}

% اضافه کردن تصویر با شرح. برای اضافه کردن تصویر بدون شرح، فقط دستور درج‌تصویر را استفاده کنید.
%\شروع{شکل}[h!]
%\تنظیم‌ازوسط
%  \درج‌تصویر[پهنا=0.5\پهنای‌متن]{img.jpg}
%\شرح{شرح تصویر}
%\برچسب{تصویر:۱}
%\پایان{شکل}
%\پایان{زیرمسائل}

%کد زیر را ببینید.
%%\lstinputlisting{source_filename.py}
%\begin{latin}
%\begin{lstlisting}
%# Import the modules
%import sys
%import random
%
%ans = True
%
%while ans:
%question = raw_input("Ask the magic 8 ball a question: (press enter to quit) ")
%
%answers = random.randint(1,8)
%
%if question == "":
%sys.exit()
%
%elif answers == 1:
%print "It is certain"
%
%elif answers == 2:
%print "Outlook good"
%
%elif answers == 3:
%print "You may rely on it"
%
%elif answers == 4:
%print "Ask again later"
%
%elif answers == 5:
%print "Concentrate and ask again"
%
%elif answers == 6:
%print "Reply hazy, try again"
%
%elif answers == 7:
%print "My reply is no"
%
%elif answers == 8:
%print "My sources say no"
%\end{lstlisting}
%\end{latin}